\section*{Розділ II. Структура та склад ЦВКс}
\addcontentsline{toc}{section}{Розділ II. Структура та склад ЦВКс}

\subsection*{2.1. Склад ЦВКс}
\addcontentsline{toc}{subsection}{2.1. Склад ЦВКс}
    2.1.1. ЦВКс складається з 5-9 членів, які обираються КСУ з числа студентів Університету.

    2.1.2. Конкретна кількість членів ЦВКс визначається КСУ при її формуванні з урахуванням обсягу завдань та кількості виборів, що плануються.

    2.1.3. Члени ЦВКс виконують свої обов'язки на громадських засадах.

\subsection*{2.2. Вимоги до кандидатів у члени ЦВКс}
\addcontentsline{toc}{subsection}{2.2. Вимоги до кандидатів у члени ЦВКс}
    2.2.1. Кандидатом у члени ЦВКс може бути студент Університету денної форми навчання, який:

        \begin{enumerate}[label=\alph*)]
            \item Не перебуває в академічній відпустці;
            \item Дав письмову згоду на участь у роботі ЦВКс.
        \end{enumerate}

    2.2.2. Кандидат у члени ЦВКс не може одночасно бути кандидатом на інші виборні посади в органах студентського самоврядування.

    2.2.3. Член ЦВКс не може одночасно обіймати посади в інших органах студентського самоврядування, крім участі в роботі комісій та робочих груп на громадських засадах.

\subsection*{2.3. Порядок формування ЦВКс}
\addcontentsline{toc}{subsection}{2.3. Порядок формування ЦВКс}
    2.3.1. ЦВКс формується КСУ шляхом обрання членів із числа кандидатів, висунутих відповідно до цього Положення.

    2.3.2. Право висування кандидатів у члени ЦВКс мають:

        \begin{enumerate}[label=\alph*)]
            \item Делегати КСУ;
            \item Органи студентського самоврядування (СР КАІ, СР СМ, СУД, СРФ/СРІ);
            \item Студентські організації, зареєстровані в Університеті;
            \item Ініціативні групи студентів у складі не менше 10 осіб.
        \end{enumerate}

    2.3.3. Кандидатури подаються до КСУ не пізніше ніж за 7 календарних днів до засідання КСУ, на якому планується обрання ЦВКс.

    2.3.4. Обрання членів ЦВКс здійснюється таємним голосуванням більшістю голосів від загального складу делегатів КСУ.

\subsection*{2.4. Структура керівництва ЦВКс}
\addcontentsline{toc}{subsection}{2.4. Структура керівництва ЦВКс}
    2.4.1. ЦВКс очолює Голова ЦВКс, який обирається членами ЦВКс зі свого складу на першому засіданні після формування комісії.

    2.4.2. З числа членів ЦВКс також обираються Секретар ЦВКс, відповідальний за ведення документообігу та протоколювання засідань, та заступник Голови ЦВКс.

    2.4.3. Голова, Секретар та заступник Голови ЦВКс обираються простою більшістю голосів від загального складу членів ЦВКс.

    2.4.4. Голова, Секретар та заступник Голови ЦВКс можуть бути переобрані на нових виборах у разі незадовільного виконання своїх обов'язків за рішенням більшості членів ЦВКс.

\subsection*{2.5. Припинення членства в ЦВКс}
\addcontentsline{toc}{subsection}{2.5. Припинення членства в ЦВКс}
    2.5.1. Членство в ЦВКс припиняється у разі:

        \begin{enumerate}[label=\alph*)]
            \item Закінчення строку повноважень ЦВКс;
            \item Подання особистої заяви про складення повноважень;
            \item Відрахування з Університету або переведення на заочну форму навчання;
            \item Порушення вимог, встановлених пунктом 2.2.1 цього Положення;
            \item Систематичного невиконання обов'язків члена ЦВКс;
            \item Рішення КСУ про дострокове припинення повноважень за поданням ЦВКс або СУД.
        \end{enumerate}

    2.5.2. У разі припинення повноважень члена ЦВКс, КСУ може обрати нового члена на строк до завершення повноважень ЦВКс даного скликання.

    2.5.3. ЦВКс вважається правомочною при наявності не менше 3 членів. У разі, якщо кількість членів ЦВКс стає менше 3, КСУ зобов'язана протягом одного місяця поповнити її склад або сформувати нову ЦВКс.